% =============================================================================
% ch09_conclusion.tex — Chapter 9: Conclusion | Biomimetic Navigator
% =============================================================================

\selectlanguage{hebrew}
\section{סיכום, מגבלות ועתיד}
\label{sec:conclusion}

\subsection{תרומות המחקר}
בעבודה זו הצגנו את \en{NavigatorCrew}, מערכת מרובת-סוכנים לניתוח ומימוש ניווט ביומימטי לרחפנים. הוכחנו כי אימוץ עקרונות אקולוקציה מעולם העטלפים, בשילוב עם אלגוריתמי \en{SLAM} מודרניים והיתוך חיישנים רב-מודאלי, מאפשר דיוק גבוה וצריכת אנרגיה נמוכה בסביבות מורכבות.

\subsection{מגבלות}
למרות התוצאות המבטיחות בסימולציה, קיימות מספר מגבלות:
\begin{itemize}
    \item \textbf{רעש סביבתי:} השפעת רעשי רוח ומנועים על חיישני הסונאר דורשת מחקר נוסף בסינון רעשים אקטיבי.
    \item \textbf{חומרת זמן-אמת:} הטמעת המודלים של \en{Vision-AI} דורשת משאבי חישוב משמעותיים (GPU) שאינם זמינים תמיד ברחפנים זעירים.
\end{itemize}

\subsection{כיוונים עתידיים}
אנו מתכננים להרחיב את המחקר לניווט נחילים (\en{Swarm Navigation}) בהשראת להקות עטלפים, ולשלב מעבדים נוירומורפיים (\en{Neuromorphic Chips}) לשיפור יעילות האנרגיה.

\subsection{מסקנות סופיות}
השילוב בין בינה מלאכותית, רובוטיקה ועקרונות ביולוגיים פותח אופקים חדשים ליצירת מכונות אוטונומיות חכמות וחסינות יותר. פלטפורמת ה-\en{NavigatorCrew} מהווה צעד משמעותי לקראת מימוש חזון זה.
